% A real number line with a red "knife" at the supremum M and a green
% knife at M - 1/2 epsilon, with scattered Cauchy-sequence terms
% a_1,...,a_8 clustering just left of M.
% Source: book/tex/calculus/limits.tex (§ Completeness / Cauchy sequences).
\documentclass[border=2mm]{standalone}
\usepackage{tikz}
\usepackage{amssymb}
\usepackage{amsmath}
\begin{document}
\begin{tikzpicture}[scale=0.7]
  \draw[<->] (-8,0)--(8,0);
  \node[below] at (8,0) {$\mathbb{R}$};
  \draw[red] (3.1,1)--(3.1,-1);
  \draw[green!45!black] (2.7,1.5)--(2.7,-1.5);
  \node[red, below right] at (3.1,-1) {$M$};
  \node[green!45!black, above] at (2.7,1.5) {$M-\tfrac12\varepsilon$};
  \foreach \x/\lbl in {-5/{a_1}, -1/{a_2}, 6/{a_3}, 0/{a_4}, 4.8/{a_5},
                       4.1/{a_6}, 3.5/{a_8}, 1.5/{a_7}} {
    \fill[blue] (\x,0) circle (2.5pt);
    \node[blue, below] at (\x,0) {$\lbl$};
  }
  \foreach \x in {3.3, 3.2, 3.05} \fill[blue] (\x,0) circle (2.5pt);
\end{tikzpicture}
\end{document}
